\documentclass[a4paper]{paper}

\usepackage[catalan]{babel}
\usepackage[
backend=biber,
citestyle=numeric  
]{biblatex}
\usepackage{fontspec}
\usepackage[margin=2cm]{geometry}
\usepackage{graphicx}
\usepackage{titlesec}
\usepackage{url}

\setlength{\parindent}{0pt}
\setlength{\parskip}{0.5em}
\def\arraystretch{1.5}

\nocite{*}
\bibliography{bibliografia.bib}

\newcommand{\sectionbreak}{\clearpage}

\begin{document}

\begin{titlepage}
	\centering
	\vspace{1cm}
	\includegraphics[width=0.25\textwidth]{images/etseib}
	\par\vspace{1cm}
	\textsc{ \LARGE Escola Tècnica Superior d'Enginyeria \\[1em] 
		Industrial de Barcelona}
	\par\vspace{2cm}
	\textbf{\LARGE Projecte II}
	\par\vspace{2cm}
	{\Huge Construcció d'una orquestra amb Raspberry Pi}
	\vfill
	\begin{flushright}
		\large
		Arnau Canyadell i Miquel \par
		Joan Marcè i Igual \par
	\end{flushright}
\end{titlepage}

\tableofcontents

\newpage

\section{Introducció}

\subsection{Motivacions}

Donat que els dos creadors d'aquest projecte som músics i que tots dos hem participat en diverses orquestres el projecte de construir una orquestra mitjançant la tecnologia i \emph{Internet of Things} ens va semblar molt interessant i que podria integrar-se bé amb els nostres coneixements.

\subsection{Projectes similars}

Buscant si hi ha altres projectes semblants al que es vol realitzar s'ha trobat que Google l'any 2011 va crear el projecte \emph{YouTube Symphony Orchestra 2011}\cite{YoutubeSymphoni} que tenia com a objectiu poder fer una orquestra a nivell mundial. Per fer-ho tots els músics es connectaven a internet i després tocaven tots junts i tot seguit totes les pistes d'àudio que aquests generaven s'ajuntaven formant una orquestra.

A part també hi ha un projecte anomenat \emph{Orchestra! a Distributed Platform for Virtual Musical Groups and Music Distance Learning over the Internet in Java \texttrademark Technology} \cite{Orchestra}. Que permet a diversos grups de música assajar o aprendre conjuntament a través d'internet tot i que no es trobin junts físicament.

\section{Objectius}
A l'inici del projecte es van establir uns objectius a complir en el projecte. Aquests es van dividir en dos grups. D'una banda, els objectius a curt termini, objectius que estaven ben definits i s'havien de fer sí o sí abans que s'acabés el projecte; i d'altra banda hi havia els objectius o idees a llarg termini, que no estaven tan ben definits i que s'intentarien completar si hi havia temps.

\subsection{Objectius a curt termini}
\begin{itemize}
	\item Aconseguir rebre la música d'un ordinador
	\item Enviar la música a una Raspberry Pi
	\item Reproduir la música en una Raspberry Pi
\end{itemize}

\subsection{Objectius a llarg termini}
\begin{itemize}
	\item Programar un bot de Telegram que reprodueixi amb l'orquestra de Raspberry Pis la música que se li enviï.
	\item Enviar música del mòbil a l'ordinador.
\end{itemize}

\section{Desenvolupament del treball}
El treball s'ha dut a terme durant un el quadrimestre de primavera, des de mitjans de febrer fins a finals de maig. S'ha estructurat la feina en 14 sessions, fent defenses parcials del projecte a les sessions 5 i 9, i una defensa final en l'última.

\subsection{Eines de treball}
Degut a ala naturalesa de desenvolupament de programari informàtic del projecte i de la necessitat de mantenir un historial de versions i un ordre del codi, s'ha treballat amb un \textbf{repositori públic de GitHub} (\url{https://github.com/jmigual/projecte2/}).

Per a l'assignació de tasques dintre del grup s'ha utilitzat l'eina Trello\cite{trello}, que permet crear tasques amb data límit, assignar-les a una persona i marcar tasques com a realitzades o endarrerides, entre altres coses.

Per al desenvolupament del programari s'ha utilitzat majoritàriament el llenguatge Python\cite{python}, amb el qual està fet tot el codi de l'orquestra, inclosos director, músics i bot. S'ha utilitzat l'entorn de Pycharm\cite{pycharm} per a escriure codi, executar-lo i depurar-lo en aquest llenguatge. Per a la configuració de les raspberry pis s'ha utilitzat shell script, i també s'ha treballat molt amb comandes de git.

Finalment, totes les presentacions del projecte i la memòria s'han fet amb \LaTeX.

\subsection{Evolució del projecte}
En el repositori hi ha un document, \texttt{documentacio.md}, on es detalla l'historial de feina feta en cada sessió (començant en la sessió 3).

\subsubsection{Inici (sessions 1 a 3)}
En les primeres sessions es va decidir el tema i es van establir els objectius del projecte. També es van definir les eines de treball que s'utilitzarien durant el projecte: es va crear el repositori a GitHub i es va crear el Trello. D'altra banda, es va buscar documentació sobre el FluidSynth i el Polyphone i es va analitzar el codi del director que ja estava implementat. Finalment, es va aprendre a configurar la raspberry pi i es va configurar la raspberry que s'ha utilitzat per fer les proves al llarg de tot el desenvolupament.

\subsubsection{Programació del director i els músics (sessions 4 a 9)}
Un cop endinsats en el projecte, es va començar a desenvolupar el programari. Es va programar, primer de tot, una versió del músic, que rebia la música del director i la reproduïa. En aquest punt es va trobar un problema amb la configuració de la sortida d'àudio de la raspberry pi: la sortida per defecte és el minijack, i no se sabia com canviar-la a la targeta USB. Per tal de solucionar-ho es va provar de canviar diverses configuracions, i finalment es va optar per instal·lar pixel a l'ordinador monoplaca i canviar la configuració manualment des de l'entorn gràfic. Així es va aconseguir solucionar el problema però el canvi implicava connectar-se manualment a l'entorn gràfic cada cop que s'iniciava la raspberry pi.

Un cop tot funcionava es modificar el director. Es van crear les classes Director i Nota per separar, ordenar i simplificar el codi. En la segona defensa del projecte, a la sessió 9, es va poder dur a terme la primera prova de l'orquestra amb diversos músics. Al final d'aquesta fase del treball ja s'havien complert els objectius a curt termini.

\subsubsection{Implementació del bot de Telegram i reproducció d'arxius MIDI (sessions 10 a 14)}
Finalment es va trobar una manera de canviar la sortida d'àudio de les raspberry pis i es va crear un script per canviar-lo de forma simplificada. Després d'això es va treballar en l'objectiu de crear un bot de Telegram que reproduís la música que se li envia, i així complir els dos objectius a llarg termini. Per a dur a terme aquest objectius es va paral·lelitzar la feina en la creació del bot i en la reproducció d'arxius MIDI.

\subsubsection{Creació del bot de Telegram}


\subsubsection{Reproducció d'arxius MIDI}


\section{Resultats}



\printbibliography

\end{document}